\item \model{مرور جامع و مبانی تحلیلی تنظیم ابرپارامترها (\en{Literature Review \& Theoretical Foundations})}

\paragraph*{مقدمه و سیر تاریخی: از تنظیم اکتشافی تا نظریه تحلیلی مقیاس‌پذیری}
آموزش مدل‌های زبانی پیشرفته (\en{LLMs}) مستلزم صرف بودجه‌های محاسباتی فراتر از $10^{25}$ تا $10^{26}$ فلاپس (\en{FLOPs}) است \cite{shulgin2026lmo}. در چنین رژیمی از مقیاس، رویکردهای سنتی مبتنی بر جستجوی شبکه‌ای (\en{Grid Search}) یا آزمون و خطا از نظر اقتصادی و زمانی ناممکن هستند. قوانین مقیاس‌پذیری اولیه نظیر فرمول‌بندی کاپلان و همکاران \cite{kaplan2020scaling} و قانون داده‌محور چینچیلا (\en{Chinchilla}) \cite{hoffmann2022chinchilla}، اگرچه رابطه توانی میان خطای آزمون و مقیاس مدل ($N$) و داده ($D$) را آشکار ساختند، اما فرایند آموزش را به عنوان یک جعبه سیاه با ابرپارامترهای بهینه‌سازی ثابت فرض می‌کردند. 

در سال‌های اخیر، پارادایم علمی حاکم بر مهندسی مدل‌های زبانی دچار یک دگرگونی بنیادین شده است: پژوهشگران با پیوند زدن معادلات دیفرانسیل تصادفی (\en{SDEs})، تحلیل طیفی ماتریس هسیان، اوراکل‌های حداقل‌سازی خطی (\en{LMO})، و نظریه تانسورپروگرام‌ها (\en{$\mu$P})، تنظیم ابرپارامترها را به یک دانش ریاضی تحلیلی، پیش‌بین و قابل انتقال (\en{Transferable}) در ابعاد مختلف بدل کرده‌اند. نمودار زیر چارچوب مفهومی انتقال همه‌جانبه ابرپارامترها را بر حسب محورهای گوناگون مقیاس به تصویر می‌کشد:

\begin{center}
\begin{tikzpicture}[
    scale=0.85, every node/.style={transform shape},
    box/.style={rectangle, rounded corners=6pt, draw=#1!80!black, fill=#1!12, thick, align=center, inner sep=7pt, font=\small\bfseries},
    bridge/.style={rectangle, rounded corners=4pt, draw=orange!80!black, fill=orange!15, dashed, thick, align=center, inner sep=5pt, font=\footnotesize},
    arrow/.style={->, >={Stealth[scale=1.1]}, thick, draw=blue!70!black},
    darrow/.style={<->, >={Stealth[scale=1.1]}, thick, draw=purple!80!black}
]
    % Nodes
    \node[box=blue, text width=3.4cm] (dense_small) at (0, 0) {مدل پایه چگال کوچک\\(\en{Small Dense Proxy})\\$d_{\text{in}}, d_{\text{out}}, L_{\text{base}}, B_0, D_0$};
    \node[box=cyan, text width=3.4cm] (dense_large) at (6.5, 0) {مدل چگال بزرگ مقیاس\\(\en{Large Dense Target})\\$\rho_d \cdot d, \rho_L \cdot L, \rho_B \cdot B, \rho_D \cdot D$};
    \node[box=teal, text width=3.4cm] (moe_dense) at (0, -4.2) {مدل متخصصان چگال\\(\en{Dense MoE Companion})\\$a = N, H = Nh, r_N = N$};
    \node[box=magenta, text width=3.4cm] (moe_sparse) at (6.5, -4.2) {مدل متخصصان تنک هدف\\(\en{Sparse MoE Target})\\$a < N, H_a = ah, r_a = a$};

    % Bridges and Arrows
    \draw[arrow] (dense_small) -- node[above, font=\footnotesize\bfseries, text=blue!80!black] {انتقال عرض و عمق (\en{Complete$^{(d)}$P})} node[below, font=\scriptsize, text=gray!80!black] {$\eta \propto d_{\text{in}}^{-1} L^{\alpha-1}$, $\lambda \propto d_{\text{in}}^{-1}$} (dense_large);
    
    \draw[darrow] (dense_small) -- node[left, font=\footnotesize\bfseries, text=purple!80!black, text width=2.4cm, align=center] {پل اول (\en{Bridge I})\\عرض فعال $\mu$P} (moe_dense);
    
    \draw[darrow] (moe_dense) -- node[above, font=\footnotesize\bfseries, text=purple!80!black] {پل دوم (\en{Bridge II}): تحلیل \en{SDE}} node[below, font=\scriptsize, text=gray!80!black] {تطابق بار داده و خنثی‌سازی خطای مرتبه اول} (moe_sparse);

    \draw[arrow, dashed, draw=red!70!black] (dense_small) to[bend left=25] node[right=4pt, font=\footnotesize\bfseries, text=red!80!black, text width=3.2cm, align=center] {قاعده طلایی \en{Complete-muE}:\\تنظیم یکباره چگال،\\انتقال مستقیم به \en{MoE}} (moe_sparse);

    \node[bridge, text width=7.2cm] at (3.25, -2.1) {ماتریس تعادل \en{Seesaw} و قوانین گام (\en{Step Law}):\\تحدب سطح اتلاف، ناوردایی $B^*$ به $N$، و هم‌ارزی $\alpha\sqrt{\beta}=\mathrm{const}$};
\end{tikzpicture}
\end{center}

\paragraph*{هندسه چشم‌انداز اتلاف و برازش تجربی: قانون گام (\en{Step Law})}
در پژوهش جامع لی و همکاران (\en{Step Law}) \cite{li2025steplaw}، طی آموزش بیش از ۳,۷۰۰ مدل زبانی چگال و تنک بر بستر ۱۰۰ تریلیون توکن با مصرف نزدیک به یک میلیون ساعت پردازش \en{H800}، ویژگی‌های ساختاری کانتورهای بهینه‌سازی اثبات گردید. فرض کنید $\mathcal{L}_{A,\mathbb{D},N,D}(\text{LR}, \text{BS})$ تابع زیان مدل زبانی با معماری $A$ و توزیع داده $\mathbb{D}$ تحت تعداد پارامترهای غیرجای‌گذاری $N$ و حجم داده $D$ باشد. مقادیر بهینه نرخ یادگیری اوج $\eta$ و اندازه دسته $B$ به صورت زیر تعریف می‌شوند:
\begin{equation}
    \eta^*, B^* = \arg\min_{\text{LR}, \text{BS}} \mathcal{L}_{A,\mathbb{D},N,D}(\text{LR}, \text{BS})
\end{equation}
یافته‌های این پژوهش موارد کلیدی زیر را اثبات نمود \cite{li2025steplaw}:
\begin{itemize}
    \item \textbf{تحدب و ناحیه فلات بهینه (\en{Loss Landscape Convexity}):} چشم‌انداز اتلاف حول جفت بهینه $(\ln \eta^*, \ln B^*)$ شدیداً محدب، کاسه‌ای شکل و دارای یک حاشیه خطای امن (\en{Plateau}) است؛ به‌طوری‌که تلورانس‌های اندک در انتخاب ابرپارامتر مانع از افت کارایی مدل نمی‌شوند.
    \item \textbf{استقلال اندازه دسته بهینه از ظرفیت مدل ($N$):} بر پایه آزمون تحلیل واریانس سلسله‌مراتبی $F$، برای برازش اندازه دسته در فضای لگاریتمی:
    \begin{align}
        \text{مدل متکی بر } N: \quad & \ln B = \beta_0 + \beta_1 \ln N \implies R^2_{\text{adj}} = -0.032 \\
        \text{مدل متکی بر } D: \quad & \ln B = \beta_0 + \beta_2 \ln D \implies R^2_{\text{adj}} = 0.821 \\
        \text{مدل دومتغیره کامل: } \quad & \ln B = \beta_0 + \beta_1 \ln N + \beta_2 \ln D \implies R^2_{\text{adj}} = 0.823
    \end{align}
    آمار آزمون برای متغیر $\ln N$ در مدل کامل برابر با $F = 0.08$ و $p = 0.257$ به دست آمد (کاملاً بی‌معنا)، در حالی که برای $\ln D$ مقدار $t = 11.863$ با سطح معناداری $p < 0.001$ تثبیت گردید.
    \item \textbf{روابط توانی قانون گام:} معادلات بهینه چندمتغیره به فرم بسته زیر حاصل شدند:
    \begin{equation}
        \eta(N, D) = 1.79 \, N^{-0.713} \, D^{0.307}, \qquad B(D) = 0.58 \, D^{0.571}
    \end{equation}
    \item \textbf{رفع سوگیری چولگی به چپ با تثبیت $\eta_{\min}$:} در زمان‌بندی‌های سنتی که نرخ پایانی به شکل کسری از نرخ اوج است ($\eta_{\min} = 0.1 \eta_{\max}$)، افزایش نرخ اوج منجر به بزرگ ماندن مخرب نرخ پایانی و عدم همگرایی در انتهای آموزش می‌گردد. تثبیت نرخ پایانی در یک مقدار بسیار کوچک (مانند $\eta_{\min} = 10^{-5}$) پیش‌شرط لازم برای حذف این سوگیری سیستماتیک است.
\end{itemize}

\paragraph*{سطح اتلاف تعاملی و غیرقابل تفکیک \en{Farseer}}
لی و همکاران در بخش دوم پژوهش‌های مقیاس‌پذیری خود تحت عنوان \en{Farseer} \cite{li2025farseer}، فرض پایه‌ای قانون چینچیلا را به چالش کشیدند. در قانون چینچیلا اتلاف به صورت تفکیک‌پذیر جمعی مدل می‌شود ($\mathcal{L} = E + A N^{-\alpha} + B D^{-\beta}$) که مستلزم صفر بودن مشتق دوم متقاطع است:
\begin{equation}
    \frac{\partial^2 \mathcal{L}_{\text{Chinchilla}}}{\partial N \partial D} = 0
\end{equation}
با آموزش ۱,۰۰۰ مدل و استفاده از سنجه نظریه اطلاعات بیت بر کاراکتر (\en{Bits Per Character - BPC}):
\begin{equation}
    \text{BPC} = \frac{1}{M} \sum_{i=1}^M \log_2 \frac{1}{P(x_i \mid x_{<i})} = \frac{\text{Loss}}{\ln 2} \times \frac{N_{\text{tokens}}}{N_{\text{chars}}}
\end{equation}
تحلیل تفاضل متناهی گسسته روی شبکه لگاریتمی $\lambda = \sqrt{2}$ تعریف شد:
\begin{equation}
    \Delta_D \mathcal{L}(N, D) = \mathcal{L}(N, D) - \mathcal{L}(N, \lambda D) \approx [V(D) - V(\lambda D)] + [H(N, D) - H(N, \lambda D)]
\end{equation}
نتایج نشان داد که $\Delta_D \mathcal{L}$ همبستگی قوی و منظمی با $N$ دارد ($R^2 = 0.9807$) که خطای برون‌یابی چینچیلا تا سقف ۴۳۳٪ را اثبات نمود. با بهره‌گیری از الگوریتم برازش تکه‌تکه‌ای تفاضلی (\en{Differential Piecewise Fitting}) در فضای توابع نمایی-کشیده (\en{Stretched-Exponential})، قانون نهایی \en{Farseer} به دست آمد:
\begin{equation}
    \mathcal{L}(N, D) = \exp\left(a_3 N^\gamma + b_3\right) + \exp\left(a_2 N^\beta + b_2\right) \cdot D^{-\exp\left(a_1 N^\alpha + b_1\right)}
\end{equation}
با مقادیر عددی برازش‌شده بر بستر معماری بهینه‌شده \en{LLaMA}:
\begin{equation}
    \mathcal{L}(N, D) = e^{-0.021 N^{0.169} - 0.091} + e^{88.01 N^{-0.1} - 6.287} \cdot D^{-e^{-0.124 N^{0.123} + 0.424}}
\end{equation}

\noindent
\textbf{اثبات یکنوایی فیزیکی (\en{Monotonicity}):} مشتق نسبت به داده‌ها همواره اکیداً منفی است:
\begin{equation}
    \frac{\partial \mathcal{L}}{\partial D} = - A(N) B(N) D^{-A(N) - 1} < 0, \quad \forall N, D > 0
\end{equation}
زیرا ضرایب $A(N) = \exp(a_1 N^\alpha + b_1)$ و $B(N) = \exp(a_2 N^\beta + b_2)$ همواره مثبت‌اند. بررسی عددی $\partial \mathcal{L}/\partial N$ نیز منفی بودن پایدار آن را در تمام بازه‌های پارامتری تأیید می‌کند. مهم‌تر از همه، قانون \en{Farseer} نشان می‌دهد که نسبت بهینه داده به پارامتر ($D/N$) برخلاف ادعای چینچیلا ثابت نیست، بلکه با رشد بودجه محاسباتی $C = 6ND$ روندی پیوسته صعودی دارد. همچنین حذف پارامترهای لایه جای‌گذاری (\en{Embedding}) از شمارش $N$، دقت برون‌یابی را تا ۴ برابر بهبود بخشید \cite{li2025farseer}.

\paragraph*{پویایی‌های زمانی و قانون چندتوانی (\en{Multi-Power Law - MPL})}
لو و همکاران \cite{luo2025multipower} مسئله پیش‌بینی پیوسته منحنی اتلاف $\mathcal{L}(t)$ در طول زمان را تحت هر برنامه زمانی دلخواه از نرخ یادگیری نزولی $E = \{\eta_1, \eta_2, \dots, \eta_T\}$ حل نمودند. 

\noindent
\textbf{اصل تطابق مجموع نرخ یادگیری (\en{LR Sum Matching}):} با نگاشت به‌روزرسانی گسسته \en{SGD} به جریان پیوسته گرادیان (\en{Gradient Flow}):
\begin{equation}
    \frac{d\theta(\tau)}{d\tau} = -\nabla \mathcal{L}(\theta(\tau)) \implies \theta_t \approx \theta\left(S_1(t)\right), \quad S_1(t) := \sum_{\tau=1}^t \eta_\tau
\end{equation}
منحنی اتلاف به دو ترم تجزیه می‌شود: اتلاف فرآیند مرجع با نرخ ثابت در گام معادل $Z(t) = S_1(t)/\eta_0$ منهای ترم افت اتلاف حاصل از کاهش نوسانات تصادفی ($LD(t)$):
\begin{equation}
    \mathcal{L}(t) = \mathcal{L}_{\text{const}}(Z(t)) - LD(t)
\end{equation}
با تحلیل تجربی زمان‌بندی‌های دوسطحی (\en{Two-stage}) و بسط تلسکوپی، قانون جامع چندتوانی (\en{MPL}) استخراج شد \cite{luo2025multipower}:
\begin{equation}
    \mathcal{L}(t) = L_0 + A \cdot \left(S_1(t) + S_W\right)^{-\alpha} - \sum_{k=1}^t B (\eta_{k-1} - \eta_k) \left[ 1 - \left( C \eta_k^{-\gamma} S_k(t) + 1 \right)^{-\beta} \right]
\end{equation}
که در آن $S_k(t) = \sum_{\tau=k}^t \eta_\tau$ و $S_W = \frac{1}{2}\eta_{\max} W$.

\noindent
\textbf{پایه تئوریک بر توابع اتلاف درجه دوم و ماتریس‌های تصادفی:} لو و همکاران نشان دادند برای تابع اتلاف درجه دوم $\mathcal{L}(\theta) = \frac{1}{2}(\theta - \theta_*)^\top H (\theta - \theta_*)$ با نویز گرادیان $\mathcal{N}(0, \Sigma)$، اگر چگالی مقادیر ویژه هسیان به صورت توانی $p(\lambda) \propto \lambda^{-\nu}$ و کوواریانس نویز مشروط به صورت $\mathbb{E}[\Sigma_{ii} \mid \lambda] \propto \lambda^{-\rho} e^{-r\lambda}$ باشد، امید ریاضی اتلاف دقیقاً ساختار چندتوانی فرمول بالا را بازتولید می‌کند و توان‌های $\alpha$ و $\beta$ از روابط زیر پیروی می‌کنند:
\begin{equation}
    \alpha = 2 - \nu - \kappa, \qquad \beta = 1 - \nu - \rho
\end{equation}
با بهینه‌سازی فرمول \en{MPL} بر بستر بردار گام‌ها، ساختار بهینه \en{Warmup-Stable-Decay (WSD)} خودبه‌خود پدیدار گشت و رگرسیون نمادین نشان داد واپاشی بهینه در فاز افت، از قانون جذر-مکعب (\en{WSDSC}) پیروی می‌کند:
\begin{equation}
    \eta_t \approx \eta_{\max} (1 - \tau)^{1.5}, \quad \tau \in [0, 1]
\end{equation}
در نقطه مقابل، با استفاده از شرایط بهینگی \en{KKT} اثبات گردید که مدل‌های رقیب نظیر \en{Momentum Law} \cite{tissue2024scaling} از نظر ریاضی دچار فروپاشی (\en{Collapse}) به زمان‌بندی‌های نامطلوب پله‌ای تک‌گامه در انتهای آموزش می‌شوند.

\paragraph*{بازبینی اندازه دسته بحرانی و گرم‌کردن دسته (\en{Batch Size Warmup})}
مریل و همکاران (\en{Critical Batch Size Revisited}) \cite{merrill2025cbs} نشان دادند که رویکرد متداول تخمین اندازه دسته بحرانی بر پایه مقیاس نویز ساده مک‌کندلیش \cite{mccandlish2018empirical}:
\begin{equation}
    B_{\text{simple}} = \frac{\text{tr}(\Sigma)}{\|G\|^2}
\end{equation}
به دلیل فرضیات ناصحیح (بهینه‌ساز \en{SGD} محض و هسیان همسانگرد $H \propto I$) در مدل‌های آموزش‌دیده با \en{AdamW} به شدت نامعتبر بوده و مقدار واقعی $B^*$ را تا چند مرتبه بزرگی کمتر برآورد می‌کند. 

آن‌ها روش \textbf{آموزش شاخه‌ای محلی (\en{Local Branched Training})} را معرفی کردند که تحت فرض بازیابی موضعی (\en{Local Recovery})، اندازه دسته بحرانی موضعی $B^*_t$ را مستقیماً از چک‌پوینت‌های میانی با انشعاب شاخه‌های کوتاه $\Delta = 2\text{B tokens}$ اندازه می‌گیرد. نتایج نشان داد:
\begin{itemize}
    \item اندازه دسته بحرانی در آغاز آموزش بسیار کوچک ($B^* \approx 0$) است، در طول ۵۰ میلیارد توکن اول به شدت جهش می‌کند و سپس حول $B^* \approx 4096$ به فلات می‌رسد.
    \item این رفتار، استراتژی \textbf{گرم‌کردن اندازه دسته (\en{Batch Size Warmup})} را توجیه می‌کند:
    \begin{equation}
        \text{If } B^*_t > 2 B_t \implies B_{t+1} = 2 B_t, \quad \eta_{t+1} = \sqrt{2} \eta_t
    \end{equation}
    اعمال این استراتژی بر مدل \en{OLMo-1B} موجب \textbf{کاهش ۴۳ درصدی گام‌های متوالی گرادیان} با حفظ کامل کیفیت و عملکرد پایین‌دستی شد \cite{merrill2025cbs}.
\end{itemize}

\paragraph*{زمان‌بند الاکلنگی (\en{Seesaw}): موازنه دقیق نرخ یادگیری و اندازه دسته}
مترز و همکاران در چارچوب نظری \model{Seesaw} \cite{meterez2026seesaw}، نخستین اثبات شرایط متناهی هم‌ارزی تضعیف نرخ یادگیری و افزایش اندازه دسته را در بهینه‌ساز \en{SGD} نرمال‌شده (\en{NSGD} به عنوان پروکسی تحلیلی برای \en{Adam}) ارائه کردند:
\begin{equation}
    \theta_{t+1} = \theta_t - \eta \frac{g_t}{\sqrt{\mathbb{E}\|g_t\|^2}}
\end{equation}
تحت \textbf{رژیم تسلط واریانس (\en{Variance Dominated Regime})}:
\begin{equation}
    \mathbb{E}\|g_t\|^2 \approx \frac{\sigma^2}{B} \text{Tr}(H) \implies \theta_{t+1} = \theta_t - \tilde{\eta} \nabla \mathcal{L}(\theta_t), \quad \tilde{\eta} = \eta \frac{\sqrt{B}}{\sigma \sqrt{\text{Tr}(H)}}
\end{equation}
آن‌ها اثبات کردند (\en{Theorem 1 \& Corollary 1}) که هر دو زمان‌بند پله‌ای با ضرایب تضعیف نرخ یادگیری $\alpha_1, \alpha_2$ و ضرایب رشد دسته $\beta_1, \beta_2$، ریسک مازاد یکسانی در پایان مراحل تولید می‌کنند مشروط بر آنکه:
\begin{equation}
    \alpha_1 \sqrt{\beta_1} = \alpha_2 \sqrt{\beta_2}
\end{equation}
با احتساب مرز ناپایداری و واگرایی مجانبی که مستلزم $\alpha \ge \sqrt{\beta}$ است، تهاجمی‌ترین و سریع‌ترین برنامه زمانی بدون واگرایی در شرایط $\alpha = \sqrt{\beta}$ محقق می‌شود که ساختار الگوریتم \model{Seesaw} را تعیین می‌کند:
\begin{equation}
    \text{هر زمان زمان‌بند مرجع نرخ را با ضریب } \alpha \text{ کاهش دهد:} \quad \eta \leftarrow \frac{\eta}{\sqrt{\alpha}}, \quad B \leftarrow B \cdot \alpha
\end{equation}

\noindent
\textbf{سقف شتاب نظری تحت واپاشی کسینوسی:} در حد پیوسته، تعداد گام‌های فرآیند متناظر با زمان‌بند کسینوسی پایه برابر است با:
\begin{equation}
    T_{\text{Seesaw}} = \int_0^T \frac{\eta(t)}{\eta_0} dt = \int_0^T \cos\left(\frac{\pi t}{2T}\right) dt = \frac{2T}{\pi}
\end{equation}
که حاکی از \textbf{کاهش بیشینه ۳۶.۳ درصدی در زمان دیواری آموزش} ($1 - \frac{2}{\pi} \approx 0.363$) در شرایط برابری کامل فلاپس است \cite{meterez2026seesaw}.

\begin{center}
\begin{tikzpicture}[
    scale=0.9, every node/.style={transform shape},
    state/.style={rectangle, rounded corners=5pt, draw=blue!70!black, fill=blue!8, thick, align=center, inner sep=6pt},
    action/.style={rectangle, rounded corners=3pt, draw=red!70!black, fill=red!10, thick, align=center, inner sep=5pt, font=\small},
    trans/.style={->, >={Stealth[scale=1.1]}, thick, draw=blue!80!black}
]
    \node[state] (start) at (0, 0) {گام زمان‌بندی پایه\\کاهش نرخ با ضریب $\alpha$};
    \node[action] (seesaw) at (5.5, 0) {موازنه الگوریتم \model{Seesaw}:\\$\eta \leftarrow \eta / \sqrt{\alpha}$\\$B \leftarrow B \cdot \alpha$};
    \node[state] (result) at (11.5, 0) {ثبات مسیر بهینه‌سازی\\کاهش گام‌ها تا سقف $36.3\%$};

    \draw[trans] (start) -- node[above, font=\footnotesize, text=red!80!black] {جایگزینی} (seesaw);
    \draw[trans] (seesaw) -- node[above, font=\footnotesize, text=green!50!black] {تسلط واریانس} (result);

    \node[draw=gray, dashed, fill=gray!5, rounded corners=4pt, text width=11cm, align=center, font=\footnotesize] at (5.75, -1.8) {
        \textbf{مرز فیزیکی و شکست فرضیه:} در دسته‌های فوق‌بزرگ، نویز گرادیان محو شده و مدل به رژیم \en{NGD} تک‌بعدی میل می‌کند؛ در این حالت چرخه پایدار $\mathcal{O}(\eta h)$ رخ می‌دهد و کاهش مستقیم $\eta$ الزامی است.
    };
\end{tikzpicture}
\end{center}

\paragraph*{پدیده جبران سریع و سوئیچ دیرهنگام دسته (\en{Fast Catch-Up \& Late Switching})}
وانگ و همکاران \cite{wang2026fastcatchup} چارچوب \textbf{قوانین مقیاس تابعی (\en{Functional Scaling Laws - FSL})} را برای زمان‌بندی پویا بنیان نهادند. بر پایه تحلیل \en{SDE} رگرسیون تانسوری هسته، پویایی اتلاف پیوسته تحت زمان‌بندی دسته $b(t)$ به فرم معادلات ولترا استخراج شد:
\begin{equation}
    \mathbb{E}[\mathcal{E}(\bar{\theta}_t)] \approx \underbrace{t^{-s}}_{\text{یادگیری سیگنال}} + \underbrace{\eta \sigma^2 \int_0^t \frac{K(t - \tau)}{b(\tau)} d\tau}_{\text{انباشت نویز}}
\end{equation}
که در آن $K(t) = (t + 1)^{-(2 - 1/\beta)}$ هسته حافظه فراموشی نویز است ($s$ توان چشمه و بیانگر سختی وظیفه، و $\beta$ توان ظرفیت مدل است). تحلیل وردشی تحت قید بودجه کل داده $\int_0^T b(t) dt = D$ به نتایج شگرف زیر منتهی شد \cite{wang2026fastcatchup}:
\begin{itemize}
    \item \textbf{وظایف آسان ($s > 1 - 1/\beta$):} زمان‌بندی بهینه دسته پیوسته صعودی است و در حالت دوسطحی، باید از همان گام نخست از اندازه دسته بزرگ استفاده کرد ($P^*_D = 0$).
    \item \textbf{وظایف سخت ($s \le 1 - 1/\beta$):} زمان‌بندی بهینه یک ساختار \textbf{رشد-پایدار (\en{Stable-Growth})} به خود می‌گیرد؛ به این معنا که مدل در بیشتر طول آموزش در حداقل اندازه دسته ($B_{\min}$) باقی می‌ماند و افزایش اندازه دسته به اواخر آموزش موکول می‌شود:
    \begin{equation}
        \frac{D - P^*_D}{D} \approx D^{-\frac{1 - 1/\beta - s}{2 - 1/\beta}} = o_D(1)
    \end{equation}
    \item \textbf{اثبات پدیده جبران سریع (\en{Fast Catch-Up Effect}):} به محض جهش از دسته کوچک $B_1$ به دسته بزرگ $B_2$ در زمان $t_*$، شکاف اتلاف $G_* = \eta \sigma^2(1/B_1 - 1/B_2)$ با سرعت بسیار زیاد میرا می‌شود:
    \begin{equation}
        \mathcal{E}_{B_1 \to B_2}(t_* + \delta) - \mathcal{E}_{B_2}(t_* + \delta) \approx G_* \, \delta^{-(1 - 1/\beta)}
    \end{equation}
    زمان لازم برای جبران ($\delta_\epsilon \approx t_*^{\frac{s}{1 - 1/\beta}}$) در مقایسه با زمان تحول پس‌زمینه مدل ($t_*$)، تفکیک مقیاس زمانی خیره‌کننده‌ای دارد: $\delta_\epsilon \ll t_*$. این اصل اثبات کرد که در مدل‌های زبانی سخت، اتلاف مدل دسته بزرگ را می‌توان با آموزش روی دسته کوچک و یک جهش دیرهنگام به دسته بزرگ دقیقاً بازتولید کرد، در حالی که حجم داده مصرفی به شدت صرفه‌جویی می‌شود \cite{wang2026fastcatchup}.
\end{itemize}

\paragraph*{انتقال همه‌جانبه ماژولار با \en{Complete$^{(d)}$P}}
ملودوزنیک و همکاران \cite{mlodozeniec2025completed} چارچوب $\mu$P را در چهار بعد اصلی مدل، عمق، دسته و زمان آموزش یکپارچه ساختند:
\begin{itemize}
    \item \textbf{اصلاح لایه‌های نرمال‌سازی ترنسفورمر:} پارامترهای ضرب‌کننده لایه‌های \en{QK-Norm} بر خلاف لایه‌های عادی، در میان تمام هدهای توجه به اشتراک گذاشته می‌شوند؛ بنابراین با تکیه بر قانون قوی اعداد بزرگ (\en{SLLN})، گرادیان پس‌انتشاریافته در مقیاس $\Theta(1)$ باقی می‌ماند و پارامتر $\epsilon$ بهینه‌ساز \en{AdamW} نباید دچار تضعیف شود. در لایه جای‌گذاری ورودی، گرادیان دارای مقیاس $\Theta(1/N)$ است و $\epsilon$ باید با ضریب $\Theta(1/N)$ مقیاس شود \cite{mlodozeniec2025completed}.
    \item \textbf{قواعد انتقال اندازه دسته و افق زمانی با \en{SDE}:} تحت گسسته‌سازی فرآیند اویلر-مارویاما با نویز نوسانی، قواعد پایدارسازی برای تغییر اندازه دسته با ضریب $\kappa$ به دست آمد:
    \begin{equation}
        \eta' = \sqrt{\kappa}\eta, \quad \lambda' = \sqrt{\kappa}\lambda, \quad k' = \kappa k
    \end{equation}
    همچنین در صورت افزایش طول آموزش در اندازه دسته ثابت، نرخ یادگیری اوج باید به فرم معکوس جذر مقیاس گردد: $\eta \propto 1/\sqrt{\kappa_{\text{tokens}}}$.
    \item \textbf{بهینه‌سازی ابرپارامترهای به ازای هر ماژول (\en{Per-Module HPO}):} به دلیل نقش متفاوت لایه‌های گوناگون ترنسفورمر در جریان اطلاعات، ملودوزنیک و همکاران برای اولین بار ۷۹ ابرپارامتر مستقل را با فاکتورگیری کرونکر نوع-عمق بهینه‌سازی کردند:
    \begin{equation}
        \log \zeta_{m,\ell}(T, N, L, B) = \log \zeta^{\text{type}}_m + \log \zeta^{\text{depth}}_\ell + \log \text{SDE}(T, B) + \log \text{CP}_m(N, L)
    \end{equation}
    این رویکرد با الگوریتم جستجوی تصادفی ناحیه اعتماد (\en{TRRS}) توانست \textbf{شتاب چشمگیر ۲.۳۱ برابری در مقیاس ۵۰M و ۱.۳۲ برابری در مدل ۷.۲B} ایجاد کند \cite{mlodozeniec2025completed}.
\end{itemize}

\paragraph*{انتقال ابرپارامترها برای بهینه‌سازهای پیش‌شرطی‌شده ماتریسی}
چیو و همکاران (\en{Matrix-Preconditioned Optimizers}) \cite{qiu2026matrix} قوانین $\mu$P را برای بهینه‌سازهای ماتریسی نظیر \en{Shampoo} \cite{gupta2018shampoo}، \en{SOAP} \cite{vyas2024soap} و \en{Muon} \cite{jordan2024muon} پایه‌گذاری کردند. با تحلیل گام اول گرادیان رنک-۱ ($G = \delta x^\top$) تحت شرط $\eta Q(G) x' = \Theta(1)$:
\begin{align}
    \textbf{Shampoo:} \quad & \eta = \Theta\left(\left(\frac{d_{\text{out}}}{d_{\text{in}}}\right)^{1 - (e_L + e_R)}\right) \\
    \textbf{Muon (Newton-Schulz):} \quad & \eta = \Theta\left(\sqrt{\frac{d_{\text{out}}}{d_{\text{in}}}}\right), \quad \epsilon = \Theta\left(\sqrt{\frac{d_{\text{in}}}{d_{\text{out}}}}\right) \\
    \textbf{SOAP (Eigenbasis Adam):} \quad & \eta = \Theta\left(\frac{T_{\text{blk}}(e_L, e_R)}{d_{\text{in}}}\right), \quad T_{\text{blk}} = b_{\text{out}}^{e_L/2} b_{\text{in}}^{e_R/2}
\end{align}

\noindent
\textbf{انحرافات عرض متناهی و نرمال‌سازی طیفی:} در گام‌های پایانی آموزش، به دلیل افزایش رتبه پایا (\en{Stable Rank}) به‌روزرسانی ناشی از عمل پیش‌شرط‌سازها بر نویز تصادفی ($\text{srank}(U) \gg 1$)، نرخ‌های $\mu$P دچار شکست موضعی می‌شوند. چیو و همکاران اثبات کردند که اعمال \textbf{نرمال‌سازی طیفی برخط (\en{Online Spectral Normalization})} بر مبنای تکرار توانی (\en{Power Iteration})، با هزینه ناچیز تنها دو ضرب ماتریس-بردار در هر لایه، انحرافات عرض متناهی را به طور کامل مهار کرده و برتری ۱.۴ برابری \en{Muon} نسبت به \en{AdamW} را تا مقیاس ۱.۴B پارامتر تثبیت می‌نماید. همچنین اثبات شد کاهش وزن مستقل بهینه باید با عرض به فرم معکوس مقیاس شود ($\lambda_{\text{WD}} = \Theta(1/d_{\text{in}})$) و نسبت بهینه توکن به پارامتر در \en{Muon} ($7.4$) به دلیل بازدهی بالاتر کمتر از \en{AdamW} ($9.6$) است \cite{qiu2026matrix}.

\paragraph*{نظریه انتقال در مدل‌های ترکیبی از متخصصان (\en{Complete-muE})}
پنگ و همکاران \cite{peng2026completemue} بزرگ‌ترین چالش انتقال را حل نمودند: در مدل‌های \en{MoE}، تغییر ساختار لایه و تعداد متخصصان فعال همزمان بار کاری توکن‌های دریافتی هر متخصص را تغییر می‌دهد ($B_{\text{exp}} \approx B a/N$). چارچوب \model{Complete-muE} با یک سیستم دوپل این گره را گشود:
\begin{itemize}
    \item \textbf{پل اول (\en{Dense FFN $\leftrightarrow$ Dense MoE}):} برقراری انطباق با مدل چگال هم‌عرض فعال ($H = Nh$) با فعال‌سازی تمام متخصصان ($a = N$) و اعمال ضریب مقیاس مسیر نرمال‌شده $r_N = N$ جهت جبران میانگین‌گیری سافت‌مکس.
    \item \textbf{پل دوم (\en{Dense MoE $\leftrightarrow$ Sparse MoE}):} گذار به حالت تنک ($a < N$) از طریق تحلیل \en{SDE} سهم هر متخصص. از آنجا که نسبت تغییر اندازه دسته و افق زمانی هر متخصص برابرند ($\rho_B^{\text{exp}} = \rho_D^{\text{exp}} = a'/a$)، ضریب تصحیح اویلر-مارویاما دقیقاً خنثی می‌شود:
    \begin{equation}
        \eta' \approx \eta \sqrt{\frac{\rho_B^{\text{exp}}}{\rho_D^{\text{exp}}}} = \eta \sqrt{\frac{a'/a}{a'/a}} = \eta
    \end{equation}
    این خنثی‌سازی تحلیلی اثبات نمود که ضرایب لایه متراکم به طور پایدار مستقیماً به مدل‌های تنک قابل انتقال هستند.
    \item \textbf{فرمول عمومی بلاک‌های ترکیبی و هیبرید:} برای مدلی با شاخه‌های اشتراکی $\mathcal{D}$ و گروه‌های تنک $\mathcal{G}$، خروجی کلی لایه چنین متوازن می‌شود:
    \begin{equation}
        y(x) = \frac{d}{H_{\text{tot}}} \left[ \sum_{m \in \mathcal{D}} o_m(x) + a \sum_{g \in \mathcal{G}} \sum_{e=1}^{N_g} \pi_{g,e}(x) o_g^{(e)}(x) \right], \quad H_{\text{tot}} = \sum_{m \in \mathcal{D}} H_m + a \sum_{g \in \mathcal{G}} h_g
    \end{equation}
    این فرمول‌بندی قاعده طلایی \textbf{«یک‌بار چگال را تنظیم کن، مستقیماً به همه انواع \en{MoE} منتقل نما» (\en{Tune Dense Once, Transfer to All})} را به اثبات رساند که به شتاب همگرایی تا سقف ۵.۵ برابر در مدل‌های زبانی بدون صرف هزینه جستجوی مجدد انجامید \cite{peng2026completemue}.
\end{itemize}

\paragraph*{نظریه بهینه‌سازی مدرن اوراکل‌های خطی (\en{LMO})}
شولگین و همکاران \cite{shulgin2026lmo} با به‌کارگیری کران‌های نوین همگرایی \en{LMO} بر پایه هندسه غیراقلیدسی کووالف برای بهینه‌سازهایی نظیر \en{SignSGD} و \en{Muon}:
\begin{equation}
    \min_{1 \le k \le K} \mathbb{E}\|\nabla f(x^k)\|_* \le \frac{\Delta_0}{\eta K} + \frac{2\rho\sigma}{\alpha\sqrt{bK}} + 2\rho\sigma\sqrt{\frac{\alpha}{b}} + \frac{7L\eta}{2} + \frac{2L\eta}{\alpha}
\end{equation}
روابط تحلیلی بسته‌ای را برای کمینه‌سازی توأم خطا استخراج کردند (جدول \ref{tab:lmo_scaling_summary}). این تحلیل نشان داد بر خلاف \en{SGD} استاندارد که در آن اندازه دسته هیچ نرخ بهینه مستقلی روی افق توکن تعریف نمی‌کند، در بهینه‌سازهای نرمال‌شده وجود ترم‌های کوپلینگ نویز و تحدید محدوده، یک اندازه دسته بهینه درون‌زا با توان $b^* \propto T^{1/2}$ (تحت مومنتوم ثابت) دیکته می‌نماید.

\paragraph*{سنتز کلان و مقایسه جامع چارچوب‌های تنظیم و انتقال ابرپارامترها}
جدول \ref{tab:lmo_scaling_summary} مقایسه تحلیلی جامعی از تمامی چارچوب‌های بررسی‌شده را بر حسب ابعاد بهینه‌سازی، مبانی ریاضی، مقیاس‌های اعتبارسنجی و دستاوردهای ثبت‌شده ارائه می‌دهد.

\begin{table}[htbp]
\centering
\caption{مقایسه و سنتز جامع چارچوب‌های نظری و تجربی تنظیم ابرپارامترها در مدل‌های زبانی بزرگ}
\label{tab:lmo_scaling_summary}
\resizebox{\textwidth}{!}{%
\begin{tabular}{lccccc}
\toprule
\textbf{نام چارچوب / مقاله} & \textbf{مرجع} & \textbf{ابرپارامترهای هدف} & \textbf{مبنای ریاضی و ابزار تحلیلی} & \textbf{مقیاس مدل و داده} & \textbf{دستاورد یا شتاب گزارش‌شده} \\
\midrule
\model{Step Law} & \cite{li2025steplaw} & $\eta(N,D), B(D)$ & آزمون $F$ و اثبات تحدب چشم‌انداز زیان & ۳,۷۰۰ مدل، تا 6.5B MoE & خطای پیش‌بینی کمتر از ۰.۰۹۴٪ \\
\model{Farseer} & \cite{li2025farseer} & $\mathcal{L}(N,D)$ و بهینگی داده & تفاضل متناهی، توابع نمایی-کشیده & ۱,۰۰۰ مدل تا 25.1B & کاهش ۴۳۳٪ خطای برون‌یابی چینچیلا \\
\model{Multi-Power Law} & \cite{luo2025multipower} & منحنی پیوسته $\mathcal{L}(t)$ & تطابق مساحت نرخ، تجزیه طیفی هسیان & مدل‌های 25M تا 1B & کشف زمان‌بندی بهینه WSDSC ($\tau^{1.5}$) \\
\model{CBS Revisited} & \cite{merrill2025cbs} & اندازه دسته بحرانی $B^*_t$ & آموزش شاخه‌ای محلی و فرض بازیابی & 1B و 7B تا 2T توکن & صرفه‌جویی ۴۳٪ گام‌های آموزش \\
\model{Seesaw} & \cite{meterez2026seesaw} & موازنه $\eta$ و $B$ & هم‌ارزی ریسک NSGD با فرض تسلط واریانس & مدل‌های 150M تا 600M & کاهش ۳۶.۳٪ زمان دیواری (سقف نظری) \\
\model{Late Switching} & \cite{wang2026fastcatchup} & زمان‌بندی پویا $b(t)$ & قوانین مقیاس تابعی (FSL)، معادلات ولترا & مدل‌های Dense/MoE تا 1.1B & اثبات سبقت سوئیچ دیرهنگام در وظایف سخت \\
\model{Complete$^{(d)}$P} & \cite{mlodozeniec2025completed} & ابرپارامترهای ماژولار & فاکتورگیری کرونکر نوع-عمق و SDE & مدل‌های 50M تا 7.2B & شتاب آموزش ۲.۳۱ برابری در مقیاس پایه \\
\model{Matrix $\mu$P} & \cite{qiu2026matrix} & $\eta, \lambda_{\text{WD}}, \epsilon$ بهینه‌ساز ماتریسی & $\mu$P رنک-۱ و نرمال‌سازی طیفی برخط & مدل‌های 190M تا 1.4B & تثبیت شتاب ۱.۴ برابری Muon بر AdamW \\
\model{Complete-muE} & \cite{peng2026completemue} & انتقال Dense به انواع MoE & سیستم دوپل و خنثی‌سازی خطای مرتبه اول SDE & مدل‌های تا 6.29B کل & قاعده Tune Dense Once، شتاب تا ۵.۵× \\
\model{LMO Theory} & \cite{shulgin2026lmo} & $\eta, \alpha, b$ در افق توکن & اوراکل حداقل‌سازی خطی و کران‌های تحلیلی & مدل‌های ترنسفورمر 160M & استخراج رابطه هذلولوی گام و اندازه دسته \\
\bottomrule
\end{tabular}%
}
\end{table}